\chapter[Conclusão]{Conclusão}
\label{conclusao}

Este trabalho investigou o problema de otimização de rotas para patrulhamento com drones, modelado como uma variante do Problema do Caixeiro Viajante com custo dependente da sequência de visita. Quatro métodos foram implementados e comparados: busca exaustiva (referência exata), Algoritmo Genético, Otimização por Enxame de Partículas e Otimização por Colônia de Formigas. Os experimentos totalizaram 5029 execuções em 30 instâncias entre 10 e 100 pontos.

\section{Principais Contribuições}

Os resultados confirmam a hipótese principal: as meta-heurísticas produzem rotas de baixo \textit{makespan} com custo computacional muito inferior ao da enumeração completa. O GA destaca-se como o método mais eficaz: encontrou o ótimo global em todas as 18 instâncias pequenas (10--15 pontos) no seu melhor caso, com desvio padrão máximo de 0,44, e manteve vantagem consistente nas instâncias maiores, com \textit{makespan} até 22\% inferior ao do ACO na instância 100a.

O PSO ocupou a segunda posição na maioria das instâncias, mas apresentou maior variabilidade entre sementes, refletindo a dependência da representação por chaves aleatórias. O ACO apresentou os maiores \textit{gaps} em relação ao ótimo, com convergência prematura nas primeiras 20--30 iterações. Esse resultado sugere que o feromônio tridimensional sofre diluição no espaço de busca, reduzindo a eficácia da construção probabilística.

A hipótese secundária também se confirmou: GA e ACO, por operarem diretamente sobre permutações, exploram a estrutura combinatória do problema de forma mais natural que o PSO. No entanto, o ACO não se beneficiou dessa vantagem devido à dispersão do feromônio no tensor tridimensional.

As contribuições específicas incluem: (i) a formalização do modelo com custo dependente da sequência e penalização angular; (ii) a implementação comparativa de quatro métodos sobre a mesma função objetivo; (iii) a base experimental com 5029 execuções e saídas estruturadas; (iv) a análise empírica detalhada de convergência, estabilidade e escalabilidade.

\section{Limitações}

Três limitações metodológicas devem ser consideradas. Primeiro, a consistência entre a documentação do tensor $G[i][j][k]$ e a ordem dos índices na geração dos custos precisa ser validada antes de atribuir significado físico definitivo à penalização de curva. Embora isso não afete a comparação entre métodos, limita a interpretação do modelo como simulação de voo real. Segundo, os hiperparâmetros foram fixados sem ajuste fino (população~100, 100 iterações para todos). O ACO, em particular, poderia se beneficiar de variações em $\alpha$, $\beta$, $\rho$ e $Q$. Terceiro, as instâncias são sintéticas com coordenadas aleatórias, o que não reflete necessariamente a geometria de cenários reais de patrulhamento.

\section{Trabalhos Futuros}

Diversas direções de continuidade são possíveis. O ajuste fino de hiperparâmetros por instância poderia alterar o ordenamento relativo entre os métodos. A implementação de variantes mais avançadas de cada meta-heurística, como o MMAS para ACO (com limites $\tau_{\min}$ e $\tau_{\max}$) ou o operador EAX para GA, pode melhorar a qualidade das soluções. A hibridização entre métodos, como no esquema GA-ACO em dois estágios \cite{hga2024hybrid}, representa uma direção promissora.

No horizonte mais amplo, a integração de aprendizado profundo com meta-heurísticas swarm -- como DeepACO \cite{deepaco2023}, NeuFACO \cite{neufaco2025} e PGACO/PPOACO \cite{ppaco2024} -- constitui o estado-da-arte em otimização combinatória neural. A aplicação dessas abordagens ao modelo com custo dependente da sequência é uma extensão natural. Por fim, a validação do modelo em cenários reais de patrulha, com dados geográficos e restrições operacionais de VANTs, permitiria avaliar o impacto prático das soluções encontradas.

\section{Contribuições em Produção Bibliográfica}

Este trabalho gerou uma base de código aberta para experimentos com meta-heurísticas em problemas de roteamento, documentada no repositório associado a esta monografia.
